% ---------------------------------------------------------------------------
% Williamson et al. (1992) Test Case 5 -- zonal flow over an isolated mountain.
% Initial condition / setup block. Requires \usepackage{amsmath}.
% ---------------------------------------------------------------------------
The Williamson et al.\ (1992) test case~5 prescribes a steady zonal (solid-body
rotation) flow impinging on an isolated mountain. The flow has no analytic
solution: the orography continually generates gravity waves and a downstream
Rossby wave train, so the dynamics remain active throughout the integration,
which makes it the canonical benchmark for long-time conservation. Let
$\mathbf{x}=(x,y,z)$ be Cartesian coordinates on the sphere of radius~$a$, with
longitude $\lambda=\operatorname{atan2}(y,x)$ and latitude
$\varphi=\arcsin(z/a)$.

The initial velocity is a solid-body rotation about the polar axis,
%
\begin{equation}
  \mathbf{u}_0 = \frac{u_0}{a}\,\bigl(-y,\;x,\;0\bigr)
              \;=\; u_0\cos\varphi\,\hat{\boldsymbol{\lambda}},
  \label{eq:w5-u}
\end{equation}
%
i.e.\ an eastward zonal jet of peak speed $u_0$. The free-surface elevation in
gradient-wind (geostrophic) balance with \eqref{eq:w5-u} is
%
\begin{equation}
  \eta_0(\varphi) = -\,\frac{1}{g}\left(a\,\Omega\,u_0 + \tfrac12 u_0^2\right)\sin^2\varphi ,
  \label{eq:w5-eta}
\end{equation}
%
and the isolated conical mountain (orography) is
%
\begin{equation}
  b(\lambda,\varphi) = h_s\left(1 - \frac{r}{R}\right),
  \qquad
  r = \min\!\Bigl(R,\;\sqrt{(\varphi-\varphi_c)^2 + (\lambda-\lambda_c)^2}\Bigr),
  \label{eq:w5-mtn}
\end{equation}
%
a cone of height $h_s$ and base radius $R$ centred at $(\lambda_c,\varphi_c)$
(zero outside $r=R$). The layer depth is then
%
\begin{equation}
  D_0 = H + \eta_0 - b ,
  \label{eq:w5-D}
\end{equation}
%
with $H$ the mean fluid depth. The parameters are
$a = 6{,}371{,}220~\mathrm{m}$,
$\Omega = 7.292\times10^{-5}~\mathrm{s^{-1}}$,
$g = 9.8~\mathrm{m\,s^{-2}}$,
$u_0 = 20~\mathrm{m\,s^{-1}}$,
$H = 5960~\mathrm{m}$,
$h_s = 2000~\mathrm{m}$,
$R = \pi/9$,
and mountain centre $(\lambda_c,\varphi_c) = (-\pi/2,\,\pi/6)$
($90^\circ$W, $30^\circ$N).
